% !TEX program = xelatex
\documentclass[12pt,a4paper]{article}

\usepackage[UTF8]{ctex}
\usepackage{amsmath}
\usepackage{geometry}
\geometry{left=2.5cm, right=2.5cm, top=2.5cm, bottom=2.5cm}
\usepackage{hyperref}
\hypersetup{colorlinks=true, linkcolor=blue, citecolor=blue, urlcolor=blue}
\usepackage{enumitem}
\usepackage{booktabs}
\usepackage{array}
\usepackage{float}
\usepackage{longtable}
\usepackage{titlesec}

\titleformat{\section}{\Large\bfseries}{\thesection}{1em}{}
\titleformat{\subsection}{\large\bfseries}{\thesubsection}{1em}{}
\titleformat{\subsubsection}{\normalsize\bfseries}{\thesubsubsection}{1em}{}

\newcommand{\recordlink}[2]{\href{#1}{#2}}

\begin{document}

\begin{center}
    {\LARGE\bfseries 基于测向机示向度的无线电干扰源自动定位与清除策略}\\[1em]
    {\large AI工具使用说明}
\end{center}

\vspace{1em}

本说明根据项目留存的AI交互记录编写，用于说明AI在论文整理、表达校核、图表辅助说明及技术概念澄清中的辅助作用。其中五份来自DeepSeek网页端，一份来自OpenCode共享页面；其中的建议仅作为讨论与编辑参考，模型构建、公式推导、程序运行、数值结果判定及论文定稿均须经团队成员独立核对后采用。

% ============================================================
\section{交互记录与工具信息}

\begin{enumerate}
    \item \textbf{DeepSeek网页端}（\url{https://chat.deepseek.com}）
    \begin{itemize}
        \item 记录模型：DeepSeek-V4.1-Flash。
        \item 记录范围：本说明列示五份AI交互记录，链接见表\ref{tab:ai-usage}。
        \item 使用内容：模型准备的公共定义整理、LaTeX表格排版与报错定位、HTTP接口概念澄清、摘要详略比较、TikZ流程图与图像说明文档的文字起草，以及公式环境与标签的排版整理。
        \item 使用方式：将题目背景、已有文字或待处理的LaTeX片段提供给AI，获得候选表述、格式建议或排错思路；随后由团队成员结合本题约束、程序输出和论文上下文进行修改与取舍。
    \end{itemize}
    \item \textbf{OpenCode共享页面}（\url{https://opncd.ai}）
    \begin{itemize}
        \item 记录模型：GLM-5.3-Flash（共享页显示）。
        \item 使用内容：将问题二中可能位置区域 $F_1$ 的陈列公式改写为带编号、可交叉引用的\texttt{equation}环境，并补充公式标签。
    \end{itemize}
\end{enumerate}

% ============================================================
\section{具体使用目的与环节}

\renewcommand{\arraystretch}{1.35}
\setlength{\LTleft}{0pt}
\setlength{\LTright}{0pt}
\begin{longtable}{@{}c >{\raggedright\arraybackslash}p{5.6cm} >{\raggedright\arraybackslash}p{3.2cm} >{\raggedright\arraybackslash}p{5.0cm}@{}}
\caption{六份代表性交互记录的用途概览}\label{tab:ai-usage}\\
\toprule
\textbf{序号} & \textbf{交互主题与辅助内容} & \textbf{对应环节} & \textbf{网址} \\
\midrule
\endfirsthead
\caption[]{六份交互记录的用途概览（续）}\\
\toprule
\textbf{序号} & \textbf{交互主题与辅助内容} & \textbf{对应环节} & \textbf{网址} \\
\midrule
\endhead
\bottomrule
\endfoot
\bottomrule
\endlastfoot
1 & 讨论是否设置“模型准备”一节，并起草坐标系、示向角域、集合算子和时间约定等共用定义；同时调整数值实例坐标表的双栏布局。 & 公共定义与论文排版 & \url{https://chat.deepseek.com/share/okjd6rmebcxs06iqfy} \\
2 & 定位六列表格最后一行多出对齐符导致的编译错误，给出与\texttt{booktabs}风格一致的修正写法。 & LaTeX表格校核 & \url{https://chat.deepseek.com/share/a4wbjd9lf39akz9iif} \\
3 & 澄清HTTP接口、浏览器客户端及跨域访问的基本概念，为理解模拟器或服务接口的调用方式提供技术背景。 & 技术概念澄清 & \url{https://chat.deepseek.com/share/p2qw3qtb76p4496dx9} \\
4 & 比较问题一、问题二两版摘要的详略差异，检查算法表述、数值验证、筛选阈值与唯一性规则是否被遗漏。 & 摘要修订与表达校核 & \url{https://chat.deepseek.com/share/xppxhsbxkvl4gax389} \\
5 & 起草顶点提取TikZ流程图，提出可独立引用的短命题，并为问题一、问题二验证图编写面向文本理解的\texttt{README.md}说明。 & 图表说明与主体格式美化 & \url{https://chat.deepseek.com/share/5e0eor5l1m0at3m8t8} \\
6 & 将问题二中 $F_1$ 的陈列公式改写为带编号的\texttt{equation}环境，并添加可供后文引用的公式标签。 & 问题二公式排版与编号 & \url{https://opncd.ai/share/z6SuAbcu} \\
\end{longtable}

% ============================================================
\section{关键交互内容及处理方式}

\subsection{公共定义与数值表格的结构整理}

\textbf{主要咨询内容：}针对四个问题共用的坐标、角度、定位区域和清除条件，讨论是否单设“模型准备”一节；同时将五个检测点及定位区域顶点的纵向表格改排为双栏表格。

\textbf{AI辅助要点：}建议将模型准备定位为简短的公共定义节，集中约定目标圆域、角度环形距离、示向角域、定位区域直径、最小覆盖圆与可靠清除判据，避免在后续问题中重复定义；表格采用左右并列的“对象--坐标”结构以改善版面比例。

\textbf{人工处理：}团队成员根据正文已有符号体系决定保留的定义范围，并核对公式、单位与题设条件的一致性。具体的顶点枚举、可靠接收区域构造和动态搜索策略仍分别置于相应问题中展开，不在公共定义节重复推导。

\subsection{LaTeX表格排错与版式统一}

\textbf{主要咨询内容：}对定位区域顶点坐标表的编译报错进行定位。

\textbf{AI辅助要点：}识别出表格列格式为\texttt{ccc|ccc}时，末行多写一个\texttt{\&}会造成单行单元格数量超出六列，并提示\texttt{\textbackslash toprule}、\texttt{\textbackslash midrule}和\texttt{\textbackslash bottomrule}依赖\texttt{booktabs}宏包。

\textbf{人工处理：}团队成员删除多余对齐符、统一两张坐标表的列数与表线风格，并在本地重新编译，确认表格、标题和交叉引用均能正常显示。

\subsection{摘要的详略比较与术语校核}

\textbf{主要咨询内容：}比较问题一、问题二两版摘要，判断压缩版是否遗漏了影响逻辑完整性的关键信息。

\textbf{AI辅助要点：}将两版摘要按算法刻画、理论结论、数值验证、候选筛选规则和唯一性保证逐项对比；提示简版中应谨慎压缩“顶点提取完备性”“最小覆盖圆判据”“几何不确定性指标的定义”“筛选阈值”以及“坐标词典序破同分规则”等内容。

\textbf{人工处理：}摘要最终以论文篇幅限制为准，保留能反映模型核心链条的结论，并由团队核对其中的复杂度、直径、残差和模拟结果均可由项目代码或正文推导支撑。AI给出的措辞建议不作为数值结论的来源。

\subsection{流程图、命题与图像说明的辅助起草}

\textbf{主要咨询内容：}为问题一顶点提取算法绘制紧凑流程图；将适合被后续问题引用的结论整理为短命题；为无法直接读取图像的工具补充验证图的文字说明。

\textbf{AI辅助要点：}给出“枚举不平行边界对--计算交点--检查全部半平面约束--去重与极角排序”的流程图候选布局；建议将“直径的顶点可达性”“直径圆非普遍覆盖”等已有结论用简短命题突出；并按图例、坐标、关键对象和论文用途组织\texttt{README.md}。

\textbf{人工处理：}命题只保留可由正文证明或算例支撑的结论，避免为形式美化新增未经证明的断言；流程图和图像说明中的坐标、直径及覆盖性结论均以实际求解输出和正文为准。

\subsection{HTTP接口概念澄清}

\textbf{主要咨询内容：}询问HTTP接口的组成，以及浏览器能否作为接口调用的客户端。

\textbf{AI辅助要点：}说明请求方法、URL、请求头、请求体、响应体和状态码的基本作用，并解释浏览器既可直接发起GET请求，也可通过页面脚本调用接口；跨域调用还受同源策略与CORS配置约束。

\textbf{人工处理：}该记录仅用于补充技术概念理解，不替代对模拟器接口文档、实际返回结果和程序日志的阅读与验证。

\subsection{问题二可行域公式的编号化排版}

\textbf{主要咨询内容：}将问题二中定义首次检测后可能位置区域 $F_1$ 的陈列公式，由双美元符号形式改为独立的\texttt{equation}环境。

\textbf{AI辅助要点：}建议沿用正文既有公式风格，将 $F_1$ 的集合定义置入\texttt{equation}环境，并增加\texttt{eq:F1\_projection}标签，便于后文通过交叉引用调用。

\textbf{人工处理：}团队成员核对公式内容、标签名称和引用位置是否与正文符号体系一致；公式环境调整仅改变排版与引用方式，不改变 $F_1$ 的数学定义或后续模型结论。

% ============================================================
\section{人工修改与质量保证}

\begin{enumerate}[label=\arabic*.]
    \item \textbf{模型层面：}定位区域的半平面刻画、顶点枚举完备性、区域直径计算、最小覆盖圆判据、可靠接收区域与候选点价值规则，均以题设、正文推导和代码实现为依据。AI输出仅用于辅助梳理表达和检查逻辑链条。

    \item \textbf{代码与数值层面：}正文中的定位区域顶点、直径、残差、模拟清除结果及运行时间均以本地代码运行或模拟器反馈为准。对话中出现的数值和坐标在写入论文前均应与程序输出逐项核对。

    \item \textbf{图表与排版层面：}表格列数、流程图节点、图例说明、公式编号和交叉引用均由团队成员在XeLaTeX环境中复核；发现版式或编译问题后，以实际编译结果为最终依据。

    \item \textbf{论文层面：}论文的研究路线、模型选择、结论表述与工程判断由团队成员负责。AI给出的候选文字、格式模板与技术解释均须结合题目要求审阅、修改后使用。
\end{enumerate}

\end{document}
